\documentclass[12pt, a4paper]{article}

% --- 常用宏包 ---
\usepackage[UTF8]{ctex}
\usepackage{amsmath, amsthm, amssymb, amsfonts}
\usepackage{geometry}
\usepackage{fancyhdr}
\usepackage{enumerate}
\usepackage{graphicx}
\usepackage{hyperref}
\usepackage{bm}
\usepackage{tcolorbox}

% --- 页面设置 ---
\geometry{left=2.5cm, right=2.5cm, top=2.5cm, bottom=2.5cm}
\pagestyle{fancy}
\fancyhf{}
\rhead{\today}
\lhead{近世代数作业}
\cfoot{\thepage}

% --- 环境定义 ---
\newtcolorbox{problem}[1]{
    colback=gray!5!white,
    colframe=gray!75!black,
    fonttitle=\bfseries,
    title=题目 #1
}

\newenvironment{solution}{
    \begin{proof}[\textbf{解}]
}{
    \end{proof}
}

% --- 个人信息 ---
\title{近世代数作业 10}
\author{李睿阳 (524120910059)}
\date{\today}

\begin{document}
\maketitle


\begin{problem}{1}
    使用长除法，对 $\mathbb{Z}_7[x]$ 中的 $f(x)=x^6+3x^5+4x^2-3x+2$，$g(x)=x^2+2x-3$，计算欧式除法算式中的 $q(x)$ 和 $r(x)$。
\end{problem}
\begin{solution}
    \begin{align*}
        f(x) &= x^4 g(x) + (x^5+3x^4+4x^2-3x+2) \\
        (x^5+3x^4+4x^2-3x+2) &= x^3 g(x) + (x^4+3x^3+4x^2-3x+2) \\
        (x^4+3x^3+4x^2-3x+2) &= x^2 g(x) + (x^3+7x^2-3x+2) \\
        (x^3+7x^2-3x+2) &= x g(x) + (5x^2+2) \\
        (5x^2+2) &= 5 g(x) + (-10x+17)
    \end{align*}
    故 $f(x) = (x^4+x^3+x^2+x+5)g(x) + (-10x+17)$。
    在 $\mathbb{Z}_7[x]$ 中，$-10 \equiv 4 \pmod 7$，$17 \equiv 3 \pmod 7$。
    \\
    则 $q(x) = x^4+x^3+x^2+x+5$，$r(x) = 4x+3$。
\end{solution}


\begin{problem}{2}
    已知 $\mathbb{Z}_5[x]$ 中多项式 $x^4+4$ 可以分解成一次因式乘积（素因子），求该分解。
\end{problem}
\begin{solution}
    在 $\mathbb{Z}_5[x]$ 中，$4 \equiv -1 \pmod 5$，有：
    \[ x^4+4 \equiv x^4-1 = (x^2+1)(x^2-1) \pmod 5 \]
    又 $x^2+1 \equiv x^2-4 = (x+2)(x-2) \pmod 5$，且 $x^2-1 = (x+1)(x-1)$。
    因此
    \[ x^4+4 = (x+1)(x+2)(x+3)(x+4) \pmod 5 \]
\end{solution}


\begin{problem}{3}
    设 $K$ 为域，任取 $f(x) \in K[x]$，证明：
    \begin{enumerate}[(1)]
        \item 当 $K$ 的特征为 $0$，$f'(x)=0 \iff f(x) \in K$。
        \item 当 $K$ 的特征为 $p$，$f'(x)=0 \iff f(x)=g(x^p)$，对某个 $g(x) \in K[x]$。
    \end{enumerate}
\end{problem}
\begin{solution}
    \begin{enumerate}
        \item[(1)] $\Longleftarrow$: 若 $f(x) = a_0 \in K$，由导数定义知其导数为 $0$。
        
        $\Longrightarrow$: 设 $f(x) = \sum_{i=0}^n a_i x^i$。则 $f'(x) = \sum_{i=1}^n ia_i x^{i-1} = 0$。
        这意味着对于所有 $1 \le i \le n$，都有 $ia_i = 0$。
        由于 $\text{char}(K) = 0$，对于 $i > 0$ 有 $i \cdot 1_K \neq 0$，故必有 $a_i = 0$ ($i=1,\dots,n$)。
        因此 $f(x) = a_0 \in K$。
        
        \item[(2)] $\Longleftarrow$: 若 $f(x) = g(x^p) = \sum b_i (x^p)^i = \sum b_i x^{pi}$，
        则 $f'(x) = \sum b_i (pi) x^{pi-1}$。
        在特征为 $p$ 的域中 $pi \equiv 0 \pmod p$，故 $f'(x) = 0$。
        
        $\Longrightarrow$: 设 $f(x) = \sum a_i x^i$。则 $f'(x) = \sum ia_i x^{i-1} = 0 \implies ia_i = 0$ 对所有 $i \ge 1$ 成立。
        若 $p \nmid i$，则 $i \cdot 1_K \neq 0 \implies a_i = 0$。
        故 $f(x)$ 中只有 $p$ 的倍数次幂项系数可能不为零，即 $f(x) = \sum a_{pk} x^{pk} = \sum a_{pk} (x^p)^k$。
        取 $g(y) = \sum a_{pk} y^k$，则 $f(x) = g(x^p)$。
    \end{enumerate}
\end{solution}
    

\end{document}
